% !TEX root = ../main.tex
%REVIEW-ADDR: W-reward-shape-counterfactual
% =============================================================================
% Appendix: Reward-Shape Counterfactual on Tool-Use ZVF
% -----------------------------------------------------------------------------
% A key concern raised against the main ZVF-collapse result is that the tool-use
% environment uses a near-binary reward (0.0 / 0.5 / 1.0), which is itself a
% confound: collapse could be driven by the reward shape rather than by model
% capacity or group size. This appendix documents a direct counterfactual
% sweep that manipulates only the reward shape and measures downstream ZVF.
% =============================================================================

\section{Reward-Shape Counterfactual Sweep}
\label{app:zvf_counterfactual}

\subsection{Hypothesis}
If reward-shape coarseness is the primary driver of ZVF saturation in the
tool-use environment, replacing the three-level reward (\(0.0\), \(0.5\), \(1.0\))
with a six-level shaped reward
\(\{0.0, 0.2, 0.4, 0.6, 0.8, 1.0\}\)
that grants partial credit for (i)~producing any JSON object, (ii)~naming a
tool, (iii)~naming the \emph{correct} tool, (iv)~producing a dict argument
block, and (v)~per-argument match should reduce ZVF by a measurable margin at
fixed group size and model scale. If ZVF is instead dominated by capacity or
group-size ceiling effects, the shaped reward should not substantially reduce
ZVF.

\subsection{Design}
We hold all other hyperparameters fixed and sweep
\(\{\text{Qwen3-4B-Instruct-2507}, \text{Qwen3-8B}\}
 \times \{\text{v1 binary-ish}, \text{v2 6-level}\}
 \times \{42, 43, 44\}\)
for a total of twelve GRPO runs
(LoRA rank \(32\), group size \(16\), batch size
\(64\) / \(128\) for 4B / 8B, \(50\) training steps).
The reward-version switch is gated by a single environment variable
(\texttt{TOOL\_USE\_REWARD\_VERSION}) so that the training harness, dataset,
and sampler are byte-identical between v1 and v2. ZVF is computed with the
same canonical script used elsewhere in this paper
(\texttt{scripts/zvf\_compute\_cross\_framework.py}); configs live under
\texttt{experiments/tool\_use\_zvf\_sweep/}.

\subsection{Results}
Table~\ref{tab:zvf_counterfactual} summarises per-run ZVF and pass rate.
Rows with em-dashes indicate runs that are in progress at the time of
submission; the sweep state is tracked at
\texttt{experiments/tool\_use\_zvf\_sweep/manifest.tsv}.

\begin{table}[h]
\centering
\small
\begin{tabular}{llrrrr}
\toprule
Model & Reward & Seed & Steps & ZVF (mean last 10 steps) & Pass\@1 \\
\midrule
Qwen3-4B & v1 & 42 & 50 & --- & --- \\
Qwen3-4B & v1 & 43 & 50 & --- & --- \\
Qwen3-4B & v1 & 44 & 50 & --- & --- \\
Qwen3-4B & v2 & 42 & 50 & --- & --- \\
Qwen3-4B & v2 & 43 & 50 & --- & --- \\
Qwen3-4B & v2 & 44 & 50 & --- & --- \\
\midrule
Qwen3-8B & v1 & 42 & 50 & --- & --- \\
Qwen3-8B & v1 & 43 & 50 & --- & --- \\
Qwen3-8B & v1 & 44 & 50 & --- & --- \\
Qwen3-8B & v2 & 42 & 50 & --- & --- \\
Qwen3-8B & v2 & 43 & 50 & --- & --- \\
Qwen3-8B & v2 & 44 & 50 & --- & --- \\
\bottomrule
\end{tabular}
\caption{Reward-shape counterfactual. ZVF is the fraction of GRPO groups
with within-group score variance below \(10^{-4}\), averaged over the final
ten training steps. \(3\) seeds per cell; the planned summary statistic is
paired-\(t\) on per-seed ZVF between v1 and v2 within each model. Em-dash
rows indicate runs still in progress; values will be spliced in from
\texttt{experiments/tool\_use\_zvf\_sweep/results.csv} on the next build.}
\label{tab:zvf_counterfactual}
\end{table}

\subsection{Pre-committed interpretation}
We pre-commit to the following reading of the results to avoid post-hoc
rationalisation:
\begin{itemize}
  \item \textbf{ZVF drops by \(\geq 0.15\) with v2 in both models.}
        Reward coarseness is a substantial contributor to collapse;
        the binary-ish reward is a partial confound and the
        ``capacity ceiling'' argument is weakened.
  \item \textbf{ZVF drops by \(<0.05\) under v2 in both models.}
        Collapse is not primarily reward-shape-driven; the
        capacity/group-size explanation is strengthened.
  \item \textbf{ZVF drops sharply in 4B but not in 8B (or vice versa).}
        The effect interacts with scale; we report the interaction as
        the headline result and retract the uniform ``reward-shape does
        not matter'' claim.
\end{itemize}

Reproduction instructions and raw wandb run URLs are listed in the
repository under \texttt{experiments/tool\_use\_zvf\_sweep/README.md}.
